\anonsection{ВВЕДЕНИЕ}
Язык Рефал-5 относится к семейству языков, основанных на переписывании термов:
программа задаётся набором правил сопоставления образца с объектным выражением и
построения результирующего выражения. На практике стоимость вычислений в таком
подходе часто определяется не только числом применённых правил, но и эффективностью
представления объектных выражений в памяти, особенно для операций добавления и
удаления термов с концов и операции конкатенации.

Выпускная квалификационная работа посвящена разработке компилятора базисного
подмножества Рефала-5 с представлением объектных выражений на основе
раскрученной двусторонней очереди. Исходный текст проходит лексический,
синтаксический и семантический анализ; исполнение реализовано интерпретирующим
циклом нормализации без кодогенерации, поэтому в технологической и аналитической
частях исполняющая часть системы для краткости именуется интерпретатором.
Реализация включает полную цепочку
трансляции, автоматическое тестирование и сравнение с двумя альтернативными
реализациями языка~--- классической Refal-5~PZ~\cite{refal5_home} на
языке~C и Refal-5J на платформе~JVM.

Объектом исследования является программная система, выполняющая трансляцию и
исполнение программ на базисном подмножестве Рефала-5. Пользователи системы~---
разработчики и исследователи, работающие с языком переписывания термов: преподаватели
и студенты при изучении формализма, авторы преобразователей программ при отладке
выходного кода рассахаривателей и других внешних инструментов. Целью работы является
разработка компилятора базисного подмножества Рефала-5, использующего раскрученную
двустороннюю очередь для представления объектных выражений, и подтверждение его
работоспособности автоматическими тестами и сравнительным замером производительности.

Исследовательская часть вводит базовые конструкции и терминологию Рефала-5, а затем
разбирает распространённые представления объектных выражений: классическое плоское
двусвязное, компромиссное с подвешенными скобками, массивное, векторно-списковое и дерево
конкатенаций. У каждого взвешиваются стоимости основных операций и связанные с ними
компромиссы; на этом фоне и обосновывается выбор раскрученной двусторонней очереди~---
структуры, которая амортизированно покрывает все ключевые операции.

Конструкторская часть фиксирует функциональные и нефункциональные требования к
системе, описывает её архитектуру, состав модулей и интерфейсы между ними. В
этой же части приводятся сценарии использования и определяются форматы
диагностических сообщений, что необходимо для последующей реализации.

Технологическая часть посвящена реализации: используется язык Java~25~LTS,
сборка автоматизирована средствами Apache~Maven~\cite{maven_guide}, тестирование
выполнено в JUnit~5~\cite{junit5_guide}. В этой части подробно разбираются
руководство пользователя, реализация раскрученной двусторонней очереди
и цикл нормализации с двустековой обработкой поля зрения.

В аналитической части представлены результаты функционального тестирования и
сравнительного замера производительности. Помимо автоматических тестов с
покрытием кода, выполнены два сравнительных прогона трёх реализаций
Рефала-5~--- разработанной в работе, классической Refal-5~PZ и
Refal-5J. Синтетический тест Loop/Select нагружает конкатенацию растущего
аккумулятора и проверяет, совпадает ли измеренное поведение с ожидаемой
линейной сложностью раскрученной очереди. Второй прогон~--- на реальной
программе рассахаривателя проекта \prog{refal-5-framework}~\cite{refal5_framework_github}~---
подтверждает корректность интерпретатора на содержательной нагрузке.
Полученные численные результаты сопоставлены с асимптотическими ожиданиями
и обсуждены применительно к выбору структуры данных.

Практическим результатом работы является программная система~--- компилятор
базисного подмножества Рефала-5 (Maven-проект \prog{refal5q-impl}), способный
выполнять как программы, написанные напрямую на базисном подмножестве, так и
результат работы внешнего рассахаривателя \prog{r5fw-desugar} проекта
\prog{refal-5-framework}~\cite{refal5_framework_github}, обеспечивая запуск
программ расширенного синтаксиса Рефала-5 после соответствующего шага
предварительной обработки.
